\documentclass[11pt]{article}

\usepackage[T1]{fontenc}
\usepackage{amsmath,amssymb,amsthm}
\usepackage[margin=1in]{geometry}
\usepackage[hidelinks]{hyperref}
\setlength{\emergencystretch}{3em}

\newtheorem{theorem}{Theorem}
\newtheorem{lemma}[theorem]{Lemma}
\newtheorem{corollary}[theorem]{Corollary}
\newtheorem{proposition}[theorem]{Proposition}
\newtheorem{conjecture}[theorem]{Conjecture}
\theoremstyle{remark}
\newtheorem{remark}[theorem]{Remark}

\newcommand{\R}{\mathbb R}
\newcommand{\Per}{\operatorname{Per}}

\title{Sharp Morrey Fragmentation of Diffuse Fast Type-II Layers}
\author{John Robert Lawson and Codex}
\date{Research round ending 25 July 2026}

\begin{document}
\maketitle

\begin{abstract}
This round quantifies the only energy-efficient exact \(q=4\) Type-II
geometry left after retained shrinking carriers were excluded.  A
peer-reviewed terminal local-energy estimate, combined with a
cube-localised \(H^1\)-to-\(L^{10/3}\) inequality, forces an amplitude layer
of volume radius \(R_j\) to fragment at the smaller Morrey scale
\(\rho_j=R_j^{27/26+o(1)}\).  Its perimeter and enstrophy obey
\[
 \Per(A_j)\gtrsim_\varepsilon
 R_j^{51/26+\varepsilon},
 \qquad
 \|\nabla u(t_j)\|_2^2\gtrsim_\varepsilon
 R_j^{-27/13+\varepsilon}.
\]
In particular, both quantities diverge after normalisation by their
single-carrier powers.  A smooth compactly supported divergence-free
packet family saturates the exponents under all static hypotheses used.
The exact \(q=4\) record gap nevertheless leaves every presently available
enstrophy, interface, cubic-flux, cutoff-flux, and viscous-cutoff floor
summable.  The live obstruction is therefore temporal and spectral:
one must either detect fixed energy below frequency
\(\rho_j^{-1}\), or make its ultraviolet failure time-resident.  The
argument is a conditional reduction under a putative first-blow-up
trajectory; it does not resolve a Clay alternative.
\end{abstract}

\section{Epistemic scope}

The local energy estimate used below is a specialisation of the
peer-reviewed energy-measure theorem of Leslie and Shvydkoy~\cite{LS18}.
The amplitude-layer ledger, localisation lemma, optimised scale,
perimeter and enstrophy deductions, sharp packet family, and summability
audit are repository deductions and await external mathematical review.
No claim of literature novelty is made.

The conclusions are conditional: if one smooth finite-energy
Navier--Stokes trajectory realises the energy-efficient exact \(q=4\)
first-record ledger, then its selected layer obeys the displayed
fragmentation bounds.  The packet family later in the report is
kinematic and consists of different fields; it is not a Navier--Stokes
trajectory.

\section{Setting and imported local-energy control}

Let \(u\) be a smooth divergence-free solution on
\(\R^3\times[0,T^*)\) of
\[
 \partial_tu+u\cdot\nabla u=-\nabla p+\nu\Delta u,
 \qquad
 \nabla\cdot u=0,
\]
where \(T^*<\infty\) is its first possible singular time.  Assume the
finite-energy bounds
\[
 \sup_{t<T^*}\|u(t)\|_2^2\le E_0,
 \qquad
 \int_0^{T^*}\|\nabla u(t)\|_2^2\,dt<\infty.
\]

Let \(t_j\uparrow T^*\) be exact \(q=4\) first-record times for
\[
 m_j:=\|u(t_j)\|_{L^{3,\infty}},
\]
so that
\begin{equation}
 \label{eq:q4-clock}
 m_j\asymp2^{2j},
 \qquad
 t_{j+1}-t_j\asymp\frac{2^{-11j}}j,
 \qquad
 T^*-t_j\lesssim\frac{2^{-11j}}j.
\end{equation}
Choose the near-extremising amplitude layer
\[
 A_j:=\{a_j<|u(t_j)|\le2a_j\},
 \qquad
 e_j:=\int_{A_j}|u(t_j)|^2\,dx,
 \qquad
 R_j:=|A_j|^{1/3}.
\]
In the energy-efficient branch,
\begin{equation}
 \label{eq:energy-efficient}
 0<c_0\le e_j\le E_0.
\end{equation}
The amplitude-layer theorem gives
\begin{equation}
 \label{eq:layer-scaling}
 a_j^2R_j^3\asymp e_j,
 \qquad
 a_j\asymp R_j^{-3/2},
 \qquad
 R_j\asymp m_j^{-2}\asymp2^{-4j}.
\end{equation}

The exact clock puts
\(\|u(t)\|_{L^{3,\infty}}\) in every temporal \(L^s\) below \(11/2\).
Interpolation with the energy-class \(L^2_tL^6_x\) bound and
the theorem of Leslie--Shvydkoy gives the following explicit estimate:
for every compact \(K\subset\R^3\),
\begin{equation}
 \label{eq:explicit-morrey}
 \sup_{\substack{T^*-r^{72/29}<t<T^*\\x\in K}}
 \int_{B_r(x)}|u(y,t)|^2\,dy
 \le C_Kr^{1/29}
\end{equation}
for sufficiently small \(r\).  Uniform spatial tightness of the energy
measures allows one fixed ball \(K_0\) such that
\begin{equation}
 \label{eq:compact-energy}
 \int_{A_j\cap K_0}|u(t_j)|^2\,dx\ge\frac{c_0}{2}
\end{equation}
for all sufficiently large \(j\).

\section{The local concentration inequality}

\begin{lemma}[Cube-localised \(L^{10/3}\) estimate]
\label{lem:local-10/3}
There is an absolute constant \(C\) such that, for every
\(f\in H^1(\R^3)\) and \(r>0\),
\begin{equation}
 \label{eq:local-10/3}
 \|f\|_{10/3}^{10/3}
 \le
 C\left(
 \sup_{x\in\R^3}\int_{B_{2r}(x)}|f|^2
 \right)^{2/3}
 \left(
 \|\nabla f\|_2^2+r^{-2}\|f\|_2^2
 \right).
\end{equation}
\end{lemma}

\begin{proof}
Tile \(\R^3\) by disjoint cubes \(Q\) of side \(r\).  Interpolation on one
cube gives
\[
 \|f\|_{L^{10/3}(Q)}^{10/3}
 \le
 \|f\|_{L^2(Q)}^{4/3}
 \|f\|_{L^6(Q)}^2.
\]
The scaled Sobolev inequality on \(Q\) is
\[
 \|f\|_{L^6(Q)}^2
 \le
 C\left(
 \|\nabla f\|_{L^2(Q)}^2
 +r^{-2}\|f\|_{L^2(Q)}^2
 \right).
\]
Every such cube lies in a ball of radius \(2r\).  Bound the first factor
by the largest local \(L^2\) mass and sum over the disjoint cubes.
\end{proof}

\section{The optimised enstrophy surcharge}

\begin{theorem}[Morrey fragmentation forces enstrophy]
\label{thm:enstrophy}
Suppose that at \(t_j\), on every fixed compact set,
\begin{equation}
 \label{eq:beta-morrey}
 \sup_x\int_{B_r(x)}|u(y,t_j)|^2\,dy
 \le C r^\beta
\end{equation}
for all sufficiently small \(r\), where \(0<\beta<3\).  Suppose the
terminal time window of this estimate contains \(t_j\) at
\begin{equation}
 \label{eq:rho-def}
 \rho_j:=L R_j^{3/(3-\beta)}.
\end{equation}
Then one sufficiently large fixed \(L\) gives
\begin{equation}
 \label{eq:enstrophy-beta}
 \boxed{
 \|\nabla u(t_j)\|_2^2
 \ge c_{\beta,L}R_j^{-6/(3-\beta)}
 }
\end{equation}
for all sufficiently large \(j\).
\end{theorem}

\begin{proof}
Choose a fixed \(\chi\in C_c^\infty(\R^3)\) equal to one on \(K_0\), and
write
\[
 f_j:=\chi u(t_j).
\]
Equations \eqref{eq:compact-energy} and \eqref{eq:layer-scaling} imply
\begin{align}
 \|f_j\|_{10/3}^{10/3}
 &\ge
 \int_{A_j\cap K_0}|u(t_j)|^{10/3}\,dx \notag\\
 &\ge
 a_j^{4/3}
 \int_{A_j\cap K_0}|u(t_j)|^2\,dx
 \ge cR_j^{-2}.
 \label{eq:lower-10/3}
\end{align}
Only balls meeting the fixed support of \(\chi\) enter
Lemma~\ref{lem:local-10/3}.  Therefore
\[
 \sup_x\int_{B_{2\rho_j}(x)}|f_j|^2
 \le C\rho_j^\beta.
\]
The cutoff also gives
\[
 \|f_j\|_2^2\le E_0,
 \qquad
 \|\nabla f_j\|_2^2
 \le C\bigl(\|\nabla u(t_j)\|_2^2+E_0\bigr).
\]
Applying \eqref{eq:local-10/3} at \(r=\rho_j\), and then using
\eqref{eq:lower-10/3}, yields
\[
 cR_j^{-2}
 \le
 C\rho_j^{2\beta/3}
 \left(
 \|\nabla u(t_j)\|_2^2+C\rho_j^{-2}
 \right).
\]
Hence
\begin{equation}
 \label{eq:enstrophy-rearranged}
 \|\nabla u(t_j)\|_2^2
 \ge
 cR_j^{-2}\rho_j^{-2\beta/3}
 -C\rho_j^{-2}.
\end{equation}
The two terms on the right have the same \(R_j\) power:
\[
 R_j^{-2}\rho_j^{-2\beta/3}
 =L^{-2\beta/3}R_j^{-6/(3-\beta)},
\]
\[
 \rho_j^{-2}
 =L^{-2}R_j^{-6/(3-\beta)}.
\]
Since \(2>2\beta/3\), increasing the fixed \(L\) makes the first
coefficient dominate the second.  This proves
\eqref{eq:enstrophy-beta}.
\end{proof}

\begin{corollary}[Explicit \(q=4\) surcharge]
\label{cor:explicit}
For the explicit estimate \eqref{eq:explicit-morrey},
\[
 \rho_j=L R_j^{87/86},
 \qquad
 \|\nabla u(t_j)\|_2^2\gtrsim R_j^{-87/43}.
\]
In particular,
\[
 R_j^2\|\nabla u(t_j)\|_2^2
 \gtrsim R_j^{-1/43}\longrightarrow\infty.
\]
\end{corollary}

\begin{proof}
Here
\[
 \beta=\frac1{29},
 \qquad
 \frac3{3-\beta}=\frac{87}{86},
 \qquad
 \frac6{3-\beta}=\frac{87}{43}.
\]
The terminal-window condition is satisfied because
\[
 \rho_j^{72/29}\asymp R_j^{108/43},
\]
whereas \eqref{eq:q4-clock} and \eqref{eq:layer-scaling} give
\[
 \frac{T^*-t_j}{\rho_j^{72/29}}
 \lesssim
 \frac1jR_j^{\,11/4-108/43}
 =
 \frac1jR_j^{41/172}
 \longrightarrow0.
\]
\end{proof}

\section{The optimised interface surcharge}

For a measurable set \(A\), write
\(\Per(A)\in[0,\infty]\) for its De Giorgi perimeter.  If an exact
dyadic endpoint gives infinite perimeter, the next lower bound is
automatic.

\begin{theorem}[Morrey fragmentation forces interface]
\label{thm:perimeter}
Under the hypotheses of Theorem~\ref{thm:enstrophy},
\begin{equation}
 \label{eq:perimeter-beta}
 \boxed{
 \Per(A_j)
 \ge
 c_{\beta,L}
 R_j^{\,2-\beta/(3-\beta)}.
 }
\end{equation}
\end{theorem}

\begin{proof}
Tile space by cubes \(Q\) of side \(\rho_j\), and retain the finitely many
cubes meeting \(K_0\).  They all lie in one fixed larger compact set.  Put
\[
 v_Q:=|A_j\cap Q|.
\]
Since \(|u|>a_j\) on \(A_j\), the local energy estimate and
\(a_j^{-2}\lesssim R_j^3\) give
\begin{equation}
 \label{eq:cube-volume}
 v_Q
 \le a_j^{-2}\int_Q|u(t_j)|^2\,dx
 \le CR_j^3\rho_j^\beta.
\end{equation}
Conversely, \eqref{eq:compact-energy}, the upper layer amplitude
\(|u|\le2a_j\), and \(a_j^{-2}\gtrsim R_j^3\) give
\begin{equation}
 \label{eq:total-volume}
 \sum_Qv_Q
 \ge |A_j\cap K_0|
 \ge cR_j^3.
\end{equation}
The ratio of the last member of \eqref{eq:cube-volume} to
\(|Q|=\rho_j^3\) is
\[
 CR_j^3\rho_j^{\beta-3}=CL^{\beta-3}.
\]
Increase the fixed \(L\), if necessary, so that \(v_Q\le|Q|/2\) for every
retained cube.  Relative isoperimetry on each cube, followed by
\eqref{eq:cube-volume} and \eqref{eq:total-volume}, yields
\begin{align*}
 \Per(A_j)
 &\ge c\sum_Qv_Q^{2/3}\\
 &\ge
 c\bigl(R_j^3\rho_j^\beta\bigr)^{-1/3}
 \sum_Qv_Q\\
 &\ge
 cR_j^2\rho_j^{-\beta/3}
 =
 c_{\beta,L}R_j^{\,2-\beta/(3-\beta)}.
\end{align*}
\end{proof}

\begin{corollary}[Explicit \(q=4\) interface]
\label{cor:explicit-perimeter}
For \(\beta=1/29\),
\[
 \Per(A_j)\gtrsim R_j^{171/86},
 \qquad
 \frac{\Per(A_j)}{R_j^2}
 \gtrsim R_j^{-1/86}\longrightarrow\infty.
\]
\end{corollary}

\begin{remark}
The absolute perimeter in Corollary~\ref{cor:explicit-perimeter} may tend
to zero.  The theorem proves divergence only after normalisation by the
perimeter \(R_j^2\) of one radius-\(R_j\) carrier.
\end{remark}

\section{Optimisation below the unavailable endpoint}

The exact clock gives every temporal exponent \(s<11/2\).  In the
Leslie--Shvydkoy interpolation, letting \(s\uparrow11/2\) while the
energy-class interpolation parameter tends to zero makes the local
energy and terminal-window exponents approach
\begin{equation}
 \label{eq:endpoint-pair}
 \beta_*=\frac19,
 \qquad
 \alpha_*=\frac{22}{9}.
\end{equation}
No endpoint integrability claim is used.

\begin{corollary}[Optimised fragmentation powers]
\label{cor:optimized}
For every \(0<b<1/9\),
\[
 \|\nabla u(t_j)\|_2^2
 \gtrsim_b R_j^{-6/(3-b)},
\]
\[
 \Per(A_j)
 \gtrsim_b R_j^{\,2-b/(3-b)}.
\]
Equivalently, for every \(\varepsilon>0\),
\begin{equation}
 \label{eq:optimized-enstrophy}
 \boxed{
 \|\nabla u(t_j)\|_2^2
 \gtrsim_\varepsilon
 R_j^{-27/13+\varepsilon},
 }
\end{equation}
\begin{equation}
 \label{eq:optimized-perimeter}
 \boxed{
 \Per(A_j)
 \gtrsim_\varepsilon
 R_j^{51/26+\varepsilon}.
 }
\end{equation}
\end{corollary}

\begin{proof}
Given \(b<1/9\), choose interpolation parameters sufficiently close to
\eqref{eq:endpoint-pair} that the published local energy exponent exceeds
\(b\).  For \(r<1\), its estimate can be weakened to \(Cr^b\).
The corresponding terminal-window exponent \(\alpha\) can be kept
arbitrarily close to \(22/9\).  At
\[
 \rho_j=L R_j^{3/(3-b)},
\]
the window still contains \(t_j\), because the limiting comparison is
strict:
\[
 \frac{3\alpha_*}{3-\beta_*}
 =\frac{33}{13}
 <\frac{11}{4}.
\]
Apply Theorems~\ref{thm:enstrophy} and~\ref{thm:perimeter}, then let
\(b\uparrow1/9\).  The limiting exponents are
\[
 \frac6{3-1/9}=\frac{27}{13},
 \qquad
 2-\frac{1/9}{3-1/9}=\frac{51}{26}.
\]
\end{proof}

The limiting scale and equal-piece ledger are
\begin{equation}
 \label{eq:limiting-ledger}
 \rho_j=R_j^{27/26+o(1)},
 \qquad
 N_j:=\left(\frac{R_j}{\rho_j}\right)^3
 =R_j^{-3/26+o(1)},
 \qquad
 N_j^{-1}=R_j^{3/26+o(1)}.
\end{equation}
No comparable-component hypothesis is used in the theorems;
\eqref{eq:limiting-ledger} is the sharp equal-piece interpretation.

\section{Sharp divergence-free packet family}

\begin{proposition}[Static sharpness]
\label{prop:sharpness}
Fix \(0<b<3\).  There is a family
\(U_R\in C_c^\infty(\R^3;\R^3)\), \(R\downarrow0\), such that
\(\nabla\cdot U_R=0\),
\[
 \|U_R\|_2^2\asymp1,
 \qquad
 \|U_R\|_{L^{3,\infty}}\asymp R^{-1/2},
\]
and
\[
 \sup_x\int_{B_r(x)}|U_R|^2\,dx\lesssim r^b
 \qquad(0<r\le1).
\]
A fixed regular amplitude band has energy \(\asymp1\), volume radius
\(\asymp R\), and
\[
 \Per(A_R)\asymp R^{\,2-b/(3-b)},
 \qquad
 \|\nabla U_R\|_2^2\asymp R^{-6/(3-b)}.
\]
\end{proposition}

\begin{proof}
Choose a nonzero divergence-free
\(\phi\in C_c^\infty(B_1;\R^3)\), with a fixed regular amplitude band of
positive measure.  Set
\[
 N\asymp R^{-3b/(3-b)},
 \qquad
 \rho:=RN^{-1/3}\asymp R^{3/(3-b)},
 \qquad
 a:=R^{-3/2}.
\]
Place \(N\) centres \(x_k\) quasi-uniformly in a fixed unit cube, with
separation
\[
 d\asymp N^{-1/3}\gg\rho,
\]
and define
\begin{equation}
 \label{eq:packet-family}
 U_R(x):=
 a\sum_{k=1}^N
 \phi\left(\frac{x-x_k}{\rho}\right).
\end{equation}
The supports are disjoint, and divergence freedom is preserved.
Direct scaling gives
\[
 \|U_R\|_2^2
 \asymp a^2N\rho^3\asymp1,
\]
\[
 \|U_R\|_{L^{3,\infty}}
 \asymp a(N\rho^3)^{1/3}
 \asymp R^{-1/2},
\]
\[
 \|\nabla U_R\|_2^2
 \asymp a^2N\rho
 \asymp R^{-6/(3-b)}.
\]
Replicating the fixed regular amplitude band gives total volume
\(\asymp N\rho^3\asymp R^3\), fixed energy, and perimeter
\[
 \asymp N\rho^2
 \asymp R^{\,2-b/(3-b)}.
\]

It remains to check the local energy bound.  If \(r\le\rho\), boundedness
of \(\phi\) gives
\[
 \int_{B_r(x)}|U_R|^2
 \lesssim R^{-3}r^3\le r^b,
\]
because \(\rho^{3-b}=R^3\).  If \(\rho\le r\le d\), one ball meets only
\(O(1)\) packets and has energy
\[
 O(N^{-1})=O(\rho^b)\le O(r^b).
\]
If \(d\le r\le1\), quasi-uniform packing gives at most
\(O(1+Nr^3)\) packets, hence energy
\[
 O(N^{-1}+r^3)=O(r^b).
\]
This proves every assertion.
\end{proof}

\begin{remark}
Proposition~\ref{prop:sharpness} shows that the powers in
Corollary~\ref{cor:optimized} are optimal under the static inputs used.
It is not a same-trajectory Navier--Stokes construction and does not
exclude a stronger dynamical theorem.
\end{remark}

\section{Exact \texorpdfstring{\(q=4\)}{q=4} summability audit}

At the limiting power, the exact record gap satisfies
\begin{equation}
 \label{eq:gap}
 \Delta t_j:=t_{j+1}-t_j
 \asymp\frac{R_j^{11/4}}j.
\end{equation}
Even if the pointwise enstrophy floor persisted for the entire gap, its
numerical contribution would be only
\begin{equation}
 \label{eq:enstrophy-gap}
 \Delta t_jR_j^{-27/13+o(1)}
 =
 \frac1jR_j^{35/52+o(1)}.
\end{equation}
This series converges geometrically.  Equation
\eqref{eq:enstrophy-gap} is not an integrated lower bound, since no such
persistence has been proved.  Its significance is that even this
optimistic persistence would not contradict the finite dissipation budget
at the available exponent.

Likewise, the limiting perimeter floor
\[
 R_j^{51/26+o(1)}
\]
is itself summable, and Proposition~\ref{prop:sharpness} attains that
scale.  A large normalised interface is therefore not an absolute
interface-budget contradiction.

The packet family has
\[
 \|U_R\|_3^3\asymp a^3R^3\asymp R^{-3/2}.
\]
Thus the full-gap cubic-action scale is
\begin{equation}
 \label{eq:cubic-gap}
 \Delta t_j\|U_{R_j}\|_3^3
 \asymp\frac1jR_j^{5/4},
\end{equation}
and inserting one cutoff gradient at \(\rho_j\) gives
\begin{equation}
 \label{eq:flux-gap}
 \frac{\Delta t_j}{\rho_j}\|U_{R_j}\|_3^3
 \asymp\frac1jR_j^{11/52+o(1)}.
\end{equation}
For the whole-space pressure representative
\[
 p=\mathcal R_a\mathcal R_b(U_aU_b),
\]
Calder\'on--Zygmund and H\"older estimates give
\[
 \int_{\R^3}|p||U_R|
 \le
 \|p\|_{3/2}\|U_R\|_3
 \le C\|U_R\|_3^3,
\]
so pressure transport has the same dimensional cutoff scale.  The
viscous cutoff scale is
\begin{equation}
 \label{eq:viscous-gap}
 \frac{\Delta t_j}{\rho_j^2}
 \asymp\frac1jR_j^{35/52+o(1)}.
\end{equation}
All these powers are positive and summable.

There is nevertheless a favourable clock inequality.  The energy-only
low-pass estimate changes frequencies below \(\rho_j^{-1}\) on the time
scale \(\rho_j^{5/2}\), whereas
\begin{equation}
 \label{eq:lowpass-gap}
 \frac{\Delta t_j}{\rho_j^{5/2}}
 =
 \frac1jR_j^{\,11/4-135/52+o(1)}
 =
 \frac1jR_j^{2/13+o(1)}
 \longrightarrow0.
\end{equation}
What is missing is a fixed energy fraction visible below a controlled
multiple of \(\rho_j^{-1}\).  Spatial Morrey fragmentation alone does not
give this: an arbitrarily small amount of energy can carry arbitrarily
large enstrophy at frequencies far above \(\rho_j^{-1}\).

\section{Findings, conjecture, and remaining obligations}

\subsection*{Robust findings, subject to review}

\begin{enumerate}
\item The explicit local-energy exponent \(1/29\) forces
\(\rho_j=R_j^{87/86}\), perimeter
\(\gtrsim R_j^{171/86}\), and enstrophy
\(\gtrsim R_j^{-87/43}\).
\item Optimisation below the unavailable temporal endpoint improves these
to \(\rho_j=R_j^{27/26+o(1)}\), normalized perimeter
\(R_j^{-1/26+o(1)}\), and normalized enstrophy
\(R_j^{-1/13+o(1)}\).
\item A smooth compactly supported divergence-free packet family saturates
all three static powers while obeying the local energy bound at every
radius.
\item Every currently available exact \(q=4\) gap-weighted static charge
retains a positive, hence summable, power of \(R_j\).
\item The low-pass clock at \(\rho_j\) is strictly longer than the record
gap by the power \(R_j^{-2/13+o(1)}\).
\end{enumerate}

\subsection*{Things still to prove}

\begin{enumerate}
\item Prove that a fixed fraction of the selected layer energy is visible
below \(C/\rho_j\), or turn its failure into a fixed ultraviolet
\(L^2\)-energy floor rather than merely large enstrophy.
\item Give that spectral alternative same-trajectory residence,
freshness, sign, or orthogonality sufficient for a nonsummable charge.
\item Control the divergent-normalised-energy branch.
\item Treat slower Type-II schedules at or below the temporal five-power
barrier.
\item Prove regularity or construct admissible breakdown for the full
Clay data class.
\end{enumerate}

\begin{conjecture}[Morrey-scale spectral dichotomy]
For an energy-efficient exact \(q=4\) selected layer with
\(\rho_j=R_j^{27/26+o(1)}\), either a fixed fraction of its energy is
coherently visible below \(C/\rho_j\) for one fixed \(C\), or a fixed
energy fraction lies above \(K_j/\rho_j\) for some \(K_j\to\infty\) and
persists long enough to pay a nonsummable physical dissipation or
transport charge.
\end{conjecture}

No part of this conjecture is proved here.  The Clay problem remains
unsolved.

\begin{thebibliography}{9}

\bibitem{LS18}
T.~M. Leslie and R. Shvydkoy,
\emph{The Energy Measure for the Euler and Navier--Stokes Equations},
Archive for Rational Mechanics and Analysis \textbf{230} (2018),
459--492.
\href{https://doi.org/10.1007/s00205-018-1250-4}
{doi:10.1007/s00205-018-1250-4}.
Accepted source: arXiv:1705.04420v4.

\end{thebibliography}

\end{document}
